% !TEX root =./ECE444_Practice_main.tex
%
% L19 practice set (Beam Steering Lab) -- reading and defending a PHASER beam
% sweep: wrapped phase tables, peak-angle quantization, and the error budget.

\fullwidth{\textbf{Learning Objective 3.5: } I can implement beam steering on the
ADALM-PHASER and verify the steered pattern against theory.
\\[4pt]\normalfont\small Show your work. A \textbf{Documentation} line at the end
is required (write \textbf{None} if you did not collaborate).}

\begin{questions}

\question \textbf{Building the phase table for a new angle.} Your HB100 locks at
$10.525\ \ghz$ and the PHASER array has $N = 8$ elements on a $d = 14\ \text{mm}$
pitch. You plan to steer to $\theta_0 = 20\mydeg$.
\begin{parts}

	\begin{minipage}[t]{0.58\linewidth}
		\part Find the wavelength and the electrical spacing $d/\lambda$.
		\begin{solution}[1.0in]
		\eq{\lambda = \frac{c}{f} = \frac{3\times10^{8}}{10.525\times10^{9}} = 0.02850\ \m = 28.50\ \text{mm}}
		\eq{\frac{d}{\lambda} = \frac{14}{28.50} = 0.491}
		\end{solution}
	\end{minipage}\hfill%
	\begin{minipage}[t]{0.37\linewidth}
		\raggedleft \vspace{12pt}
		$\lambda$ = \ansbox[1.3in]{$28.5\ \text{mm}$}\\[10pt]
		$d/\lambda$ = \ansbox[1.3in]{$0.491$}
	\end{minipage}

	\begin{minipage}[t]{0.58\linewidth}
		\part Compute the progressive element phase $\Delta\phi$ required to
		steer the beam to $20\mydeg$.
		\begin{solution}[1.1in]
		\eq{\Delta\phi = kd\sinp{\theta_0} = 360\mydeg\ \frac{d}{\lambda}\ \sinp{\theta_0}}
		\eq{= 360\mydeg (0.491)(0.3420) = 176.8\mydeg (0.3420) = 60.5\mydeg}
		\end{solution}
	\end{minipage}\hfill%
	\begin{minipage}[t]{0.37\linewidth}
		\raggedleft \vspace{12pt}
		$\Delta\phi$ = \ansbox[1.3in]{$60.5\mydeg$}
	\end{minipage}

	\part Complete the table of phases the beamformer will apply, wrapping each
	value into $0\mydeg$ to $360\mydeg$. State which element is the first to wrap.

	\vspace{4pt}
	\noindent\begin{tabular}{l|cccccccc}
	Element & 1 & 2 & 3 & 4 & 5 & 6 & 7 & 8 \\ \hline
	Applied phase & & & & & & & & \\[10pt]
	\end{tabular}
	\vspace{2pt}

		\begin{solution}[1.0in]
		Element $n$ carries $(n-1)\Delta\phi$ with $\Delta\phi = 60.5\mydeg$, wrapped modulo $360\mydeg$:
		\eq{0.0\mydeg,\ 60.5\mydeg,\ 121.0\mydeg,\ 181.4\mydeg,\ 241.9\mydeg,\ 302.4\mydeg,\ 2.9\mydeg,\ 63.4\mydeg}
		Element 7 is the first to wrap: its unwrapped value is $6(60.5\mydeg) = 362.9\mydeg$, which the shifter produces as $2.9\mydeg$.
		\end{solution}

	\part A classmate says the wrapped table cannot produce the same beam as the
	unwrapped ramp, because element 7 now leads element 1 by only $2.9\mydeg$
	instead of lagging a full turn behind. In one or two sentences, explain why
	the two tables give identical patterns.
		\begin{solution}[1.0in]
		The field radiated by an element depends on its phase only through
		$e^{j\phi}$, which is unchanged by adding or subtracting any whole number
		of $360\mydeg$. A $362.9\mydeg$ shift and a $2.9\mydeg$ shift are the same
		complex weight, so every term in the array sum is identical and the
		patterns are identical. Wrapping matters only because the hardware cannot
		store a number larger than $360\mydeg$.
		\end{solution}

\end{parts}

\newpage

\question \textbf{Reading a measured sweep.} A cadet places the HB100 at
$+30.0\mydeg$ on the protractor arc, hunts for the peak in the Beam Steering box,
and then runs a full sweep on the Rectangular tab. The trace peaks at a commanded
angle of $28.1\mydeg$, its half-power crossings are at $20.5\mydeg$ and
$35.8\mydeg$, the peak sits $0.7\ \db$ below the frozen boresight trace, and the
first sidelobes read $-11.8$ dBc. The steer resolution is the default
$2.8125\mydeg$.
\begin{parts}

	\begin{minipage}[t]{0.58\linewidth}
		\part Find the measured half-power beamwidth and compare it with the
		value predicted from the $13.1\mydeg$ boresight beamwidth.
		\begin{solution}[1.1in]
		\eq{\theta_\text{HP,meas} = 35.8\mydeg - 20.5\mydeg = 15.3\mydeg}
		The $1/\cosp{\theta_0}$ broadening predicts
		\eq{\frac{13.1\mydeg}{\cosp{30\mydeg}} = \frac{13.1\mydeg}{0.866} = 15.1\mydeg}
		The measurement is $0.2\mydeg$ wide, well inside one grid step.
		\end{solution}
	\end{minipage}\hfill%
	\begin{minipage}[t]{0.37\linewidth}
		\raggedleft \vspace{12pt}
		$\theta_\text{HP}$ = \ansbox[1.3in]{$15.3\mydeg$}\\[10pt]
		predicted = \ansbox[1.3in]{$15.1\mydeg$}
	\end{minipage}

	\begin{minipage}[t]{0.58\linewidth}
		\part Express the difference between the physical angle and the commanded
		angle at the peak as a fraction of one sweep grid step.
		\begin{solution}[1.0in]
		\eq{\Delta\theta = 30.0\mydeg - 28.1\mydeg = 1.9\mydeg}
		\eq{\frac{1.9\mydeg}{2.8125\mydeg} = 0.68 \text{ grid steps}}
		\end{solution}
	\end{minipage}\hfill%
	\begin{minipage}[t]{0.37\linewidth}
		\raggedleft \vspace{12pt}
		offset = \ansbox[1.3in]{$0.68$ steps}
	\end{minipage}

	\part The cadet reports the array as broken because the beam did not go to
	$30\mydeg$. In two or three sentences, say whether the data support that
	conclusion.
		\begin{solution}[1.1in]
		They do not. The sweep only samples commanded angles on a $2.8125\mydeg$
		grid, so the reported peak is the grid point nearest the true peak and can
		be wrong by half a step, $1.41\mydeg$, on the grid alone. Adding the
		protractor and aiming error of roughly $1.5\mydeg$ puts a $1.9\mydeg$
		disagreement inside the expected spread. A pointing fault would show up as
		an offset of several degrees that persists at every commanded angle.
		\end{solution}

	\part The peak fell $0.7\ \db$ when the beam was steered to $30\mydeg$, while
	the projected-aperture rule predicts a drop of $10\mylog{\cosp{30\mydeg}} = -0.6\ \db$ instead.
	Give the most likely reason the measured loss is the larger of the two, and
	state whether you would expect the gap to grow or shrink at $50\mydeg$.
		\begin{solution}[1.2in]
		The projected-aperture rule is the ideal-element bound: it assumes every
		patch radiates equally in all scan directions and charges only for the
		foreshortened aperture. A real patch element has its own pattern that rolls
		off away from broadside, closer to $\cos^{1.3}\theta$ to
		$\cos^{1.5}\theta$ in power, so the measured scan loss is always at least
		as large as the bound. The gap grows with scan angle, so a $50\mydeg$ scan
		would show a larger excess than the $0.1\ \db$ seen here.
		\end{solution}

\end{parts}

\newpage

\question \textbf{Bounding the error sources.} You need to state an uncertainty
for the peak angle you measured. Work with the default $2.8125\mydeg$ steer
resolution, a source range of $1.00\ \m$, and an HB100 nominal frequency of
$10.525\ \ghz$.
\begin{parts}

	\begin{minipage}[t]{0.58\linewidth}
		\part From the sweep grid alone, give the worst-case error in the peak
		angle read off the trace, and the RMS error if the true peak is equally
		likely to fall anywhere between two grid points.
		\begin{solution}[1.1in]
		The reported peak is the nearest grid point, so the worst case is half a
		step, and a uniform spread of width $\Delta$ has standard deviation
		$\Delta/\sqrt{12}$:
		\eq{\Delta\theta_\text{max} &= \frac{2.8125\mydeg}{2} = 1.41\mydeg \\ \Delta\theta_\text{rms} &= \frac{2.8125\mydeg}{\sqrt{12}} = 0.81\mydeg}
		\end{solution}
	\end{minipage}\hfill%
	\begin{minipage}[t]{0.37\linewidth}
		\raggedleft \vspace{12pt}
		worst case = \ansbox[1.3in]{$1.41\mydeg$}\\[10pt]
		RMS = \ansbox[1.3in]{$0.81\mydeg$}
	\end{minipage}

	\begin{minipage}[t]{0.58\linewidth}
		\part The phases were computed for $10.525\ \ghz$ and applied for
		$\theta_0 = 45\mydeg$. The module then drifts to $10.325\ \ghz$. Find the
		angle the beam points at, and the resulting error.
		\begin{solution}[1.2in]
		The ramp is fixed, so the beam sits where $kd\sinp{\theta} = \Delta\phi$ at
		the new frequency:
		\eq{\theta &= \arcsin\lp \frac{f_0}{f}\sinp{\theta_0} \rp = \arcsin\lp \frac{10.525}{10.325}(0.7071) \rp \\ &= \arcsin(0.7208) = 46.1\mydeg, \qquad \Delta\theta = +1.1\mydeg}
		\end{solution}
	\end{minipage}\hfill%
	\begin{minipage}[t]{0.37\linewidth}
		\raggedleft \vspace{12pt}
		$\theta$ = \ansbox[1.3in]{$46.1\mydeg$}\\[10pt]
		$\Delta\theta$ = \ansbox[1.3in]{$+1.1\mydeg$}
	\end{minipage}

	\begin{minipage}[t]{0.58\linewidth}
		\part How far along the arc must the tripod be misplaced to introduce a
		$1.0\mydeg$ error in the reference angle?
		\begin{solution}[1.0in]
		\eq{s = R\,\theta_\text{rad} = (1.00\ \m)\lp \frac{1.0\pi}{180} \rp = 0.0175\ \m}
		which is $17\ \text{mm}$ of lateral placement error.
		\end{solution}
	\end{minipage}\hfill%
	\begin{minipage}[t]{0.37\linewidth}
		\raggedleft \vspace{12pt}
		$s$ = \ansbox[1.3in]{$17\ \text{mm}$}
	\end{minipage}

	\part Multipath in the classroom adds roughly $1\mydeg$ of peak-angle
	wander. Combine all four contributions in quadrature, then say which one you
	could reduce from the GUI and how, and which one you could not.
		\begin{solution}[1.3in]
		Taking $1.41\mydeg$ (grid), $1.5\mydeg$ (aim), $1.1\mydeg$ (drift) and
		$1.0\mydeg$ (multipath):
		\eq{\sqrt{1.41^2 + 1.5^2 + 1.1^2 + 1.0^2} = 2.5\mydeg}
		The grid term is the one the GUI controls: a smaller Steer Resolution, or
		more Phase Shift Bits, samples the trace more finely. Drift cannot be fixed
		from the GUI, since the HB100 free-runs with no reference input.
		\end{solution}

\end{parts}

\newpage

\question \textbf{Working backwards from the phase read-back.} A lab partner
recorded the Phase Control values below but forgot which steer angle was
commanded. The HB100 was at $10.525\ \ghz$.

\vspace{4pt}
\noindent\begin{tabular}{l|cccccccc}
Element & 1 & 2 & 3 & 4 & 5 & 6 & 7 & 8 \\ \hline
Applied phase & $0.0\mydeg$ & $125.0\mydeg$ & $250.1\mydeg$ & $15.1\mydeg$ & $140.1\mydeg$ & $265.2\mydeg$ & $30.2\mydeg$ & $155.2\mydeg$ \\
\end{tabular}
\vspace{6pt}

\begin{parts}

	\begin{minipage}[t]{0.58\linewidth}
		\part Recover the progressive phase $\Delta\phi$ from the table.
		\begin{solution}[1.1in]
		Take successive differences and add $360\mydeg$ wherever the value drops:
		\eq{125.0\mydeg,\ 125.1\mydeg,\ (375.1 - 250.1) = 125.0\mydeg,\ \ldots}
		Every step is the same, so
		\eq{\Delta\phi = 125.0\mydeg}
		\end{solution}
	\end{minipage}\hfill%
	\begin{minipage}[t]{0.37\linewidth}
		\raggedleft \vspace{12pt}
		$\Delta\phi$ = \ansbox[1.3in]{$125.0\mydeg$}
	\end{minipage}

	\begin{minipage}[t]{0.58\linewidth}
		\part Find the commanded steer angle.
		\begin{solution}[1.1in]
		\eq{\sinp{\theta_0} = \frac{\Delta\phi}{360\mydeg (d/\lambda)} = \frac{125.0\mydeg}{176.8\mydeg} = 0.7071}
		\eq{\theta_0 = 45.0\mydeg}
		\end{solution}
	\end{minipage}\hfill%
	\begin{minipage}[t]{0.37\linewidth}
		\raggedleft \vspace{12pt}
		$\theta_0$ = \ansbox[1.3in]{$45.0\mydeg$}
	\end{minipage}

	\part The same read-back is also consistent with $\Delta\phi = 125.0\mydeg +
	360\mydeg = 485.0\mydeg$. Explain in one or two sentences why that second
	solution can be discarded for this array.
		\begin{solution}[1.1in]
		A progressive phase of $485.0\mydeg$ would require
		$\sinp{\theta_0} = 485.0/176.8 = 2.74$, which no real angle satisfies.
		Because $d/\lambda = 0.491$ is less than one half, every wrapped phase
		table maps to exactly one beam direction in visible space, and the
		ambiguity does not arise. Spacings above a half wavelength are where the
		second solution becomes real, and that is the grating lobe of L26.
		\end{solution}

	\begin{minipage}[t]{0.58\linewidth}
		\part The partner later finds the HB100 was running at
		$10.300\ \ghz$ while these phases were applied. Where did the beam
		point?
		\begin{solution}[1.2in]
		At $10.300\ \ghz$, $\lambda = 29.13\ \text{mm}$ and $d/\lambda = 0.4807$, so
		\eq{\sinp{\theta_0} = \frac{125.0\mydeg}{360\mydeg (0.4807)} = \frac{125.0\mydeg}{173.0\mydeg} = 0.7226}
		\eq{\theta_0 = 46.3\mydeg}
		The beam sat $1.3\mydeg$ past the commanded angle.
		\end{solution}
	\end{minipage}\hfill%
	\begin{minipage}[t]{0.37\linewidth}
		\raggedleft \vspace{12pt}
		$\theta_0$ = \ansbox[1.3in]{$46.3\mydeg$}
	\end{minipage}

\end{parts}

\newpage

\question \textbf{Two frozen traces.} A group freezes the boresight sweep, then
steers to $45\mydeg$ and sweeps again. The boresight trace has a half-power width
of $13.1\mydeg$ and defines the $0\ \db$ reference. The steered trace peaks at $45.0\mydeg$,
is $18.9\mydeg$ wide at half power, and its peak is $1.7\ \db$ below the boresight
peak.
\begin{parts}

	\begin{minipage}[t]{0.58\linewidth}
		\part Predict the half-power width at $45\mydeg$ and compare it with the
		measurement.
		\begin{solution}[1.1in]
		The beam broadens as $1/\cosp{\theta_0}$, so scaling the measured $13.1\mydeg$
		broadside width gives
		\eq{\theta_\text{HP} \approx \frac{13.1\mydeg}{\cosp{45\mydeg}} = \frac{13.1\mydeg}{0.7071} = 18.5\mydeg}
		The trace reads $18.9\mydeg$, which is $0.4\mydeg$ wide and inside one grid step.
		Starting instead from the theory width $0.886\lambda/(Nd) = 13.2\mydeg$ the same
		rule gives $18.7\mydeg$, so either chain lands inside one grid step of the trace.
		\end{solution}
	\end{minipage}\hfill%
	\begin{minipage}[t]{0.37\linewidth}
		\raggedleft \vspace{12pt}
		$\theta_\text{HP}$ = \ansbox[1.3in]{$18.5\mydeg$}
	\end{minipage}

	\begin{minipage}[t]{0.58\linewidth}
		\part Predict the peak drop from the projected-aperture rule and compare
		it with the measured $1.7\ \db$ drop.
		\begin{solution}[1.0in]
		\eq{\Delta G &= 10\mylog{\cosp{45\mydeg}} \\ &= 10\mylog{0.7071} = -1.5\ \db}
		The measurement is $0.2\ \db$ deeper, consistent with the element pattern
		rolling off faster than the ideal cosine.
		\end{solution}
	\end{minipage}\hfill%
	\begin{minipage}[t]{0.37\linewidth}
		\raggedleft \vspace{12pt}
		$\Delta G$ = \ansbox[1.3in]{$-1.5\ \db$}
	\end{minipage}

	\part In one or two sentences, name the single geometric fact behind both the
	wider beam and the lower peak.
		\begin{solution}[0.9in]
		Steering tilts the beam off the face of the aperture, so the array presents
		a projected length of $Nd\cosp{\theta_0}$ rather than $Nd$ to the incoming
		wave. A shorter effective aperture gives a proportionally wider beam and
		collects proportionally less power, which is the same $\cosp{\theta_0}$
		appearing in both results.
		\end{solution}

	\begin{minipage}[t]{0.58\linewidth}
		\part A radar specification calls for a $13\mydeg$ beam maintained out to
		$45\mydeg$ of scan. How many elements would this array need at the same
		$14\ \text{mm}$ spacing, and state the cost of that choice in one phrase.
		\begin{solution}[1.2in]
		The projected aperture must equal the present one at $45\mydeg$:
		\eq{N\cosp{45\mydeg} &= 8 \\ N &= \frac{8}{0.7071} = 11.3 \rightarrow 12 \text{ elements}}
		The cost is hardware: twelve elements need twelve LNAs and phase shifters,
		three ADAR1000 chips instead of two, and a wider aperture that still loses
		the same $1.5\ \db$ of scan loss.
		\end{solution}
	\end{minipage}\hfill%
	\begin{minipage}[t]{0.37\linewidth}
		\raggedleft \vspace{12pt}
		$N$ = \ansbox[1.3in]{$12$ elements}
	\end{minipage}

\end{parts}

\vspace{1.5em}
\noindent\textbf{Documentation:} \rule[-2pt]{0.6\linewidth}{0.4pt}

\end{questions}
